\begin{enumerate}

    
    \item 已知双曲线 $\Gamma  : \dfrac{{x}^{2}}{2} - \dfrac{{y}^{2}}{2} = 1$ ,过点 $M\left( {m,0}\right)$ 作不垂直于 $x$ 轴的直线 $l$ 交双曲线 $\Gamma$ 于 $A$ 、 $B$ 两点.
    \begin{enumerate}
        \item 求双曲线离心率
        \item 若点 $A\left( {\sqrt{3},1}\right)$ ,点 $B$ 在双曲线的右支上,且 $B$ 是 ${AM}$ 的中点,求直线 $l$ 的斜率
        \item 若 $m > 0,{F}_{1},{F}_{2}$ 分别是双曲线 $\Gamma$ 的左右焦点, ${A}^{\prime }$ 是 $A$ 关于 $y$ 轴的对称点,若存在直线 $l$使得 $\vv{{F}_{1}{A}^{\prime }} \cdot  \vv{{F}_{2}B} = 0$ ,求 $m$ 的取值范围.
    \end{enumerate}
    \begin{jieda}
        \begin{enumerate}
            \item $c = 2, a = \sqrt{2}, e = \dfrac{c}{a} = \sqrt{2}$
            \item $B\left(\dfrac{\sqrt{3}+m}{2}, \dfrac{1}{2}\right)$在双曲线上，$\left(\dfrac{\sqrt{3}+m}{2}\right)^2 - \dfrac{1}{4} = 2$，故$(\sqrt{3}+m)^2 =9$又因为$B$在右支上，所以$\dfrac{\sqrt{3}+m}{2} > 0$，因此$m = 3-\sqrt{3}$，即$B\left(\dfrac{3}{2}, \dfrac{1}{2}\right)$，因此$k_l = \dfrac{1-\dfrac{1}{2}}{\sqrt{3}-\dfrac{3}{2}} = \dfrac{1}{2\sqrt{3}-3} = \dfrac{2\sqrt{3}+3}{3}$
            \item $F_1(-2, 0)$，$F_2(2, 0)$
            
            当$l$为$y = 0$时不满足要求，故设直线$l:  x = ty + m$，设$A(x_1, y_1), B(x_2, y_2)$，因此$A'(-x_1, y_1)$

            $\vv{F_1A'} = (2-x_1, y_1)$，$\vv{F_2A} = (x_2-2, y_2)$，$\vv{F_1A'} \cdot \vv{F_2A} =(2-x_1)(x_2-2)+y_1y_2 = 0$，即$(2-m-ty_1)(ty_2+m-2)+y_1y_2 = 0$，即$(1-t^2)y_1y_2+(2-m)t(y_1+y_2)-(m-2)^2 = 0$

            联立直线与双曲线$\left\{\begin{array}{l}
                 x = ty+m  \\
                 x^2-y^2 = 2
            \end{array}\right.$，消去$x$得$(ty+m)^2-y^2 = 2$，即$(t^2-1)y^2+2tmy+(m^2-2) = 0$，$t \neq \pm 1$，$\Delta = 4t^2m^2 - 4(m^2-2)(t^2-1) = 4m^2+8t^2-8 > 0$

            根据韦达定理有$\left\{\begin{array}{l}
                 y_1+y_2 = \dfrac{2tm}{1-t^2} \\
                 y_1y_2 = \dfrac{2-m^2}{1-t^2}
            \end{array}\right.$，代入$(1-t^2)y_1y_2+(2-m)t(y_1+y_2) -(m-2)^2 = 0$得$(1-t^2)\cdot\dfrac{2-m^2}{1-t^2}+(2-m)t\cdot\dfrac{2tm}{1-t^2}-(m-2)^2 = 0$，即$\dfrac{t^2-m^2+2m-1}{1-t^2} = 0$，故$t^2 = (m-1)^2$

            代入$4m^2+8t^2-8 > 0$可得$m^2+2(m-1)^2-2>0$，解得$m < 0$或$m > \dfrac{4}{3}$，又因为$m > 0, t^2 \neq 1$，所以$m \in \left(\dfrac{4}{3}, 2\right) \cup (2, +\infty)$
        \end{enumerate}
    \end{jieda}
    \item 定义 $\Gamma  : {x}^{2} - {y}^{2} = 1,{F}_{1}\text{、}{F}_{2}$ 为 $\Gamma$ 的左右焦点 
    \begin{enumerate}
        \item 求点$(2,0)$到 $\Gamma$ 渐近线的距离
        \item ${P}$ 为 $\Gamma$ 上一点, $\vv{P{F}_{1}} \cdot  \vv{P{F}_{2}} = 1$ ,求 $\bigtriangleup P{F}_{1}{F}_{2}$ 的面积
        \item 设 $\Omega  : {x}^{2} - {y}^{2} = 1$ ,其中 $\left\{  \begin{array}{l} x < 0 \\  y \leq   - 1 \end{array}\right.$ 或 $\left\{  \begin{array}{l} x > 0 \\  y \geq   - 1 \end{array}\right.$ ,过点 ${F}_{2}$ 的直线 $l$ 交 $\Omega$ 于 $P\text{、}Q$ 两点(分别位于一、四象限),过点 ${F}_{2}$ 直线 $m$ 交$\Omega$ 于 $M\text{、}N$ 两点 (分别位于三、四象限),是否存在正数 $\lambda$ ,对于任意的 $l$ ,都存在唯一的 $m$ ,使 $\left| {MN}\right|  = \lambda \left| {PQ}\right|$ 成立,若存在,求出所有的 $\lambda$ ,若不存在,请说明理由.
    \end{enumerate}
    \begin{jieda}
        \begin{enumerate}
            \item 渐近线为$y = \pm x$，$(2, 0)$到渐近线的距离为$d = \sqrt{2}$
            \item $F_1(-\sqrt{2}, 0), F_2(\sqrt{2}, 0)$
            \begin{dingautolist}{172}
                \item 设$P(x_0, y_0)$，$\vv{PF_1} \cdot \vv{PF_2} = |\vv{OP}|^2 - 2 = 1$，故$\left\{\begin{array}{l}
                     x_0^2 + y_0^2 = 3 \\
                     x_0^2 - y_0^2 = 1
                \end{array}\right.$，故$y_0 = \pm 1$
                \item 设$P\left(\dfrac{1}{cos\theta}, tan\theta\right), \theta \in \left[0, 2\pi\right)$且$\theta \neq \dfrac{\pi}{2}, \dfrac{3\pi}{2}$，则$\vv{PF_1} = \left(-\sqrt{2} - \dfrac{1}{cos\theta}, -tan\theta\right)$，$\vv{PF_2} = \left(\sqrt{2} - \dfrac{1}{cos\theta}, -tan\theta\right)$，$1 = \vv{PF_1} \cdot \vv{PF_2} = \dfrac{1}{cos^2\theta} - 2 + tan^2\theta$，即$3cos^2\theta = 1 + sin^2\theta$，即$cos^2\theta = \dfrac{1}{2}$，故$y_P = \pm 1$
            \end{dingautolist}
            故$S_{\triangle{PF_1F_2}} = \sqrt{2}$
            
            \item 由题可知$x_M < 0 < x_N$，$x_Q < \sqrt{2} < x_P$，$k_l \in (1, +\infty)$，或$k_l$不存在，$k_m \in \left[\dfrac{\sqrt{2}}{4}, 1\right)$
            
            $|F_2M| = \sqrt{(x_M-\sqrt{2})^2+y_M^2} = -\sqrt{2}x_M+1$，同理$|F_2N| = \sqrt{2}x_N - 1$，故$|MN| = 2-\sqrt{2}(x_M+x_N)$

            $|F_2P| = \sqrt{2}x_P - 1$，$|F_2Q| = \sqrt{2}x_Q - 1$，故$|PQ| = \sqrt{2}(x_P + x_Q) - 2$

            设直线$PQ: y = k_l(x-\sqrt{2})$，与双曲线联立得$(1-k_l^2)x^2+2\sqrt{2}k_l^2x-(2k_l^2+1)=0$

            $\Delta>0$, $x_P+x_Q=-\dfrac{2\sqrt{2}k_l^2}{1-k_l^2}$，故$|PQ|= \dfrac{-4k_l^2}{1-k_l^2}-2 \in (2,+\infty)$, 当$k_l$不存在时，$|PQ|=2$，故$|PQ|\in [2,+\infty)$

            同理可得$|MN|=2+\dfrac{4k_m^2}{1-k_m^2}$，故$|MN|$随$k_m$增大严增，$|MN|\in \left[\dfrac{18}{7},+\infty\right)$

            因此当$\lambda \geq \dfrac{9}{7}$满足要求.
        \end{enumerate}
    \end{jieda}
    
    \item 抛物线 $C : y^2 = 4x$ 的焦点为 $F$ ,点 $D\left(2, 0\right)$ ,过 $F$ 的直线交 $C$ 于 $A,B$ 两点. 设直线 ${AD},{BD}$ 与 $C$ 的另一个交点分别为 $M,N$.
    \begin{enumerate}
        \item 证明: 直线$MN$过定点.
        \item 设直线$MN$与直线$AB$交于点$P$，求$P$点轨迹.
    \end{enumerate}
    \begin{jieda}
        \begin{enumerate}
            \item $F(1, 0)$，$A(x_1, y_1), B(x_2, y_2), M(x_3, y_3), N(x_4, y_4)$
            
            易知直线$AB$斜率不为0，设直线$AB: x = m_1y+1$，则有$4(m_1y+1) = y^2$，即$y^2-4m_1y-4 = 0, \Delta > 0$，根据韦达定理有$y_1y_2 = -4$.
            
            易知直线$AM$斜率不为0，设直线$AM: x = m_2y+2$则有$4(m_2y+2) = y^2$，即$y^2-4m_2y-8 = 0, \Delta > 0$，根据韦达定理有$y_1y_3 = -8$，同理可得$y_2y_4 = -8$，故$y_3 = 2y_2, y_4 = 2y_1$，因此$y_3y_4 = 4y_1y_2 = -16$

            设直线$MN$的直线方程为$x = my+n$，则有$4(my+n) = y^2$，即$y^2-4my-4n = 0, \Delta > 0$，根据韦达定理有$y_3y_4 = -4n$，故$-4n = -16$，$n = 4$，因此$MN$恒过定点$(4, 0)$

            \item 直线$AB: x = m_1y+1$，直线$MN: x = my+4$
            
            $m_1 = \dfrac{x_1-x_2}{y_1-y_2} = \dfrac{\dfrac{y_1^2}{4} - \dfrac{y_2^2}{4}}{y_1-y_2} = \dfrac{y_1+y_2}{4}$，同理$m = \dfrac{y_3+y_4}{4}$，故$m = 2m_1$

            \begin{dingautolist}{172}
                \item 当$m=m_1=0$时，两直线没有交点
                \item $m_1 = \dfrac{x_P - 1}{y_P}, m = \dfrac{x_P - 4}{y_P}$，故$2\dfrac{x_P - 1}{y_P} = \dfrac{x_P - 4}{y_P}$，故$x_P = -2(y_P \neq 0)$，因此$P$的轨迹是$x = -2(y \neq 0)$.
            \end{dingautolist}
            
        \end{enumerate}
    \end{jieda}

    \item 双曲线 ${x}^{2} - \dfrac{{y}^{2}}{{b}^{2}} = 1\left( {b > 0}\right)$ 的左、右焦点分别为 ${F}_{1}$ 、 ${F}_{2}$ ,直线 $l$ 过 ${F}_{2}$ 且与双曲线交于 $A$ 、 $B$ 两点.
    \begin{enumerate}
        \item 若双曲线的离心率为 2 ；求 $b$ 的值
        \item 若 $l$ 的倾斜角为 $\dfrac{\pi }{2}$ , $\bigtriangleup {F}_{1}{AB}$ 是等边三角形,求双曲线的渐近线方程
        \item 设 $b = \sqrt{3}$ ,若 $l$ 的斜率存在,且 $\left( {\vv{{F}_{1}A} + \vv{{F}_{1}B}}\right)  \cdot  \vv{AB} = 0$ ,求  $l$ 的斜率.
    \end{enumerate}
    \begin{jieda}
        \begin{enumerate}
            \item $e = \dfrac{c}{a} = 2,a = 1,\therefore c = 2$ .$\therefore {b}^{2} = {c}^{2} - {a}^{2} = 3,\therefore b = \sqrt{3}$ 
            
            \item 设 $A\left( {{x}_{A},{y}_{A}}\right)$ . 由题意, ${F}_{2}\left( {c,0}\right) ,c = \sqrt{1 + {b}^{2}},{y}_{A}^{2} = {b}^{2}\left( {{c}^{2} - 1}\right)  = {b}^{4}$ ,因为 $\bigtriangleup {F}_{1}{AB}$ 是等边三角形,所以 ${2c} = \sqrt{3}\left| {y}_{A}\right|$ ,即 $4\left( {1 + {b}^{2}}\right)  = 3{b}^{4}$ ,解得 ${b}^{2} = 2$ . 故双曲线的渐近线方程为 $y =  \pm  \sqrt{2}x$
            \item 
             $\left( {\vv{{F}_{1}A} + \vv{{F}_{1}B}}\right)  \cdot  \vv{AB} = \left( {\vv{{F}_{1}A} + \vv{{F}_{1}B}}\right) \cdot \left( {\vv{{F}_{1}B} - \vv{{F}_{1}A}}\right) = |\vv{F_1B}|^2 - |\vv{F_1A}|^2 = 0$，故$|F_1A| = |F_1B|$

             由已知, ${F}_{1}\left( {-2,0}\right) ,{F}_{2}\left( {2,0}\right)$ .设$P(x_0, y_0)$为双曲线上一点，则$|F_1P| = \sqrt{(x_0+2)^2+y_0^2} = |1+2x_0|$

             若$A, B$位于同一支，则必有横坐标相同，与$l$ 的斜率存在矛盾，故不妨设$A$在左支，$B$在右支，故有$-(1+2x_A) = 1+2x_B$，因此$x_A+x_B = -1$.
            
            设 $A\left( {{x}_{1},{y}_{1}}\right) ,B\left( {{x}_{2},{y}_{2}}\right)$ ,直线 $l : y = k\left( {x - 2}\right)$ . 显然 $k \neq  0$ .
            
            由 $\left\{  \begin{array}{l} {x}^{2} - \dfrac{{y}^{2}}{3} = 1 \\  y = k\left( {x - 2}\right)  \end{array}\right.$ ,得 $\left( {{k}^{2} - 3}\right) {x}^{2} - 4{k}^{2}x + 4{k}^{2} + 3 = 0$ .因为 $l$ 与双曲线交于两点,所以 ${k}^{2} - 3 \neq  0$ ,且 $\Delta  = {36}\left( {1 + {k}^{2}}\right)  > 0$.

            根据韦达定理有$x_A+x_B = \dfrac{4k^2}{k^2-3} = -1$，故$k = \pm \dfrac{\sqrt{15}}{5}$
        \end{enumerate}
    \end{jieda}
    
    \item 已知平面内一动点 $P$ 到点 ${F}\left( {1,0}\right)$ 的距离与点 $P$ 到 $y$ 轴的距离的差等于 1.
    \begin{enumerate}
        \item 求动点 $P$ 的轨迹 $C$ 的方程
        \item 过点 $F$ 作两条斜率存在且互相垂直的直线 ${l}_{1},{l}_{2}$ ,设 ${l}_{1}$ 与轨迹 $C$ 相交于点 $A,B$ , ${l}_{2}$ 与轨迹 $C$ 相交于点 $D,E$ ,求 $\vv{AD} \cdot  \vv{EB}$ 的最小值.
    \end{enumerate}
    \begin{jieda}
        \begin{enumerate}
            \item 设动点 $P$ 的坐标为$(x, y)$,由题意为 $\sqrt{{\left( x - 1\right) }^{2} + {y}^{2}} - \left| x\right|  = 1$. 化简得 ${y}^{2} = {2x} + 2\left| x\right|$ ,
            
            当 $x \geq  0$ 时, ${y}^{2} = {4x}$ ；当 $x < 0$ 时, $y = 0$ .
            
            所以动点 ${P}$ 的轨迹 ${C}$ 的方程为 ${y}^{2} = {4x}\left( {x \geq  0}\right)$ 和 $y = 0\left( {x < 0}\right)$ .
            
            \item 
            设 $A\left( {{x}_{1},{y}_{1}}\right) ,B\left( {{x}_{2},{y}_{2}}\right)$ , $D\left( {{x}_{3},{y}_{3}}\right), E\left( {{x}_{4},{y}_{4}}\right)$

            
            $\vv{AD} \cdot  \vv{EB} = \left( {\vv{AF} + \vv{FD}}\right)  \cdot  \left( {\vv{EF} + \vv{FB}}\right) = \vv{AF} \cdot  \vv{EF} + \vv{AF} \cdot  \vv{FB} + \vv{FD} \cdot  \vv{EF} + \vv{FD} \cdot  \vv{FB} = \left| \vv{AF}\right|  \cdot  \left| \vv{FB}\right|  + \left| \vv{FD}\right|  \cdot  \left| \vv{EF}\right| = \left( {{x}_{1} + 1}\right) \left( {{x}_{2} + 1}\right)  + \left( {{x}_{3} + 1}\right) \left( {{x}_{4} + 1}\right)
            $
            
            由题意知,直线 ${l}_{1}$ 的斜率存在且不为 0,设为 $k$ ,则 ${l}_{1}$ 的方程为 $y = k\left( {x - 1}\right)$ .
            
            由 $\left\{  \begin{array}{l} y = k\left( {x - 1}\right) \\  {y}^{2} = {4x} \end{array}\right.$ ,得 ${k}^{2}{x}^{2} - \left( {2{k}^{2} + 4}\right) x + {k}^{2} = 0$，根据韦达定理有$\left\{\begin{array}{l}
                 {x}_{1} + {x}_{2} = 2 + \dfrac{4}{{k}^{2}}  \\
                 {x}_{1}{x}_{2} = 1
            \end{array}\right.$ 根据两直线垂直可知两直线斜率乘积为$-1$，则可得 $\left\{\begin{array}{l}
                 {x}_{3} + {x}_{4} = 2 + 4k^2  \\
                 {x}_{3}{x}_{4} = 1
            \end{array}\right.$
            
            故$\vv{AD} \cdot  \vv{EB} = 1 + \left( {2 + \dfrac{4}{{k}^{2}}}\right)  + 1 + 1 + \left( {2 + 4{k}^{2}}\right)  + 1 = 8 + 4\left( {{k}^{2} + \dfrac{1}{{k}^{2}}}\right)  \geq  8 + 4 \times  2\sqrt{{k}^{2} \cdot  \dfrac{1}{{k}^{2}}} = {16}
            $
            
            当且仅当 ${k}^{2} = \dfrac{1}{{k}^{2}}$ 即 $k =  \pm  1$ 时, $\vv{AD} \cdot  \vv{EB}$ 取最小值 16.
        \end{enumerate}
    \end{jieda}
    
    \item 已知椭圆 $C : \dfrac{{x}^{2}}{4} + \dfrac{{y}^{2}}{3} = 1$
    \begin{enumerate}
        \item 点$H(0, \sqrt{3})$，过$H$做两条相互垂直的直线，分别交椭圆于$P,Q$两点 (异于点 $H$ ),证明: $PQ$过定点.
        \item 点 $H\left( {1,\dfrac{3}{2}}\right)$ 在 $C$ 上,过点 $T\left( {4,0}\right)$ 的直线 ${l}_{2}$ 交椭圆 $C$ 于 $P,Q$ 两点 (异于点 $H$ ),设直线 ${HP}$，${HQ}$ 的斜率分别为 ${k}_{1},{k}_{2}$ ,证明: ${k}_{1} + {k}_{2}$ 为定值.
    \end{enumerate}
    \begin{jieda}
        \begin{enumerate}
            \item 根据对称性可知若$PQ$过定点，则必在$y$轴上，根据特值可知定点坐标为$M\left(0, -\dfrac{\sqrt{3}}{7}\right)$

            设$HP: y = k_1x+\sqrt{3}$，$HQ: y = k_2x+\sqrt{3}$，$k_1k_2 = -1$

            联立$HP$和椭圆可得$3x^2+4(k_1x+\sqrt{3})^2 = 12$，即$(3+4k_1^2)x^2+8\sqrt{3}k_1x = 0$

            故$x_P = \dfrac{-8\sqrt{3}k_1}{3+4k_1^2}$，故$P\left(\dfrac{-8\sqrt{3}k_1}{3+4k_1^2}, \dfrac{3\sqrt{3}-4\sqrt{3}k_1^2}{3+4k_1^2}\right)$，因此$\vv{MP} = \left(\dfrac{-8\sqrt{3}k_1}{3+4k_1^2}, \dfrac{24\sqrt{3}(1-k_1^2)}{7(3+4k_1^2)}\right)$

            同理可得$\vv{MQ} = \left(\dfrac{-8\sqrt{3}k_2}{3+4k_2^2}, \dfrac{24\sqrt{3}(1-k_2^2)}{7(3+4k_2^2)}\right) = \left(\dfrac{8\sqrt{3}k_1}{3k_1^2+4}, \dfrac{24\sqrt{3}(k_1^2-1)}{7(3k_1^2+4)}\right)$

            易知$\vv{MP} // \vv{MQ}$，因此$P,M,Q$三点共线，即$PQ$过定点为$M\left(0, -\dfrac{\sqrt{3}}{7}\right)$.
            \item 
            \begin{dingautolist}{172}
                \item 当直线 ${l}_{2}$ 的斜率为 0 时,方程为 $y = 0$ ,
                此时 $P,Q$ 两点坐标为 $\left( {-2,0}\right) ,\left( {2,0}\right)$ ,又 $H\left( {1,\dfrac{3}{2}}\right)$ ,
                则 ${k}_{1} + {k}_{2} = \dfrac{\dfrac{3}{2} - 0}{1 - \left( {-2}\right) } + \dfrac{\dfrac{3}{2} - 0}{1 - 2} =  - 1$ .
                当直线 ${l}_{2}$ 的斜率不为 0 时,由已知设直线 ${l}_{2} : x = {my} + 4\left( {m \neq  0}\right)$ ,
                设点 $P\left( {{x}_{3},{y}_{3}}\right) ,Q\left( {{x}_{4},{y}_{4}}\right)$ 且与点 $H\left( {1,\dfrac{3}{2}}\right)$ 不重合,
                联立直线 ${l}_{2}$ 与椭圆 $C$ 的方程 $\left\{  \begin{array}{l} x = {my} + 4 \\  \dfrac{{x}^{2}}{4} + \dfrac{{y}^{2}}{3} = 1 \end{array}\right.$ ,消去 $x$ ,得 $\dfrac{{\left( my + 4\right) }^{2}}{4} + \dfrac{{y}^{2}}{3} = 1$ ,
                整理得 $\left( {3{m}^{2} + 4}\right) {y}^{2} + {24my} + {36} = 0$ ,则 $\Delta  = {\left( {24}m\right) }^{2} - {144}\left( {3{m}^{2} + 4}\right)  > 0$ ,即 ${m}^{2} - 4 > 0$ ,
                解得 $m > 2$ 或 $m <  - 2$ ,且 ${y}_{3} + {y}_{4} =  - \dfrac{24m}{3{m}^{2} + 4},{y}_{3}{y}_{4} = \dfrac{36}{3{m}^{2} + 4}$ ,
                则 ${k}_{1} + {k}_{2} = \dfrac{{y}_{3} - \dfrac{3}{2}}{{x}_{3} - 1} + \dfrac{{y}_{4} - \dfrac{3}{2}}{{x}_{4} - 1} = \dfrac{{y}_{3} - \dfrac{3}{2}}{m{y}_{3} + 3} + \dfrac{{y}_{4} - \dfrac{3}{2}}{m{y}_{4} + 3}$
                $= \dfrac{\left( {{y}_{3} - \dfrac{3}{2}}\right) \left( {m{y}_{4} + 3}\right)  + \left( {{y}_{4} - \dfrac{3}{2}}\right) \left( {m{y}_{3} + 3}\right) }{\left( {m{y}_{3} + 3}\right) \left( {m{y}_{4} + 3}\right) } = \dfrac{{2m}{y}_{3}{y}_{4} + \left( {3 - \dfrac{3}{2}m}\right) \left( {{y}_{3} + {y}_{4}}\right)  - 9}{{m}^{2}{y}_{3}{y}_{4} + {3m}\left( {{y}_{3} + {y}_{4}}\right)  + 9},$
                代入 ${y}_{3} + {y}_{4} =  - \dfrac{24m}{3{m}^{2} + 4},{y}_{3}{y}_{4} = \dfrac{36}{3{m}^{2} + 4}$ ,
                得 ${k}_{1} + {k}_{2} = \dfrac{{2m} \times  \dfrac{36}{3{m}^{2} + 4} - \left( {3 - \dfrac{3}{2}m}\right)  \times  \dfrac{24m}{3{m}^{2} + 4} - 9}{{m}^{2} \times  \dfrac{36}{3{m}^{2} + 4} - {3m} \times  \dfrac{24m}{3{m}^{2} + 4} + 9} = \dfrac{\dfrac{9\left( {{m}^{2} - 4}\right) }{3{m}^{2} + 4}}{\dfrac{9\left( {{m}^{2} - 4}\right) }{3{m}^{2} + 4}} =  - 1$ .
                综上, ${k}_{1} + {k}_{2}$ 为定值,且 ${k}_{1} + {k}_{2} =  - 1$ .
                \item 以$H$为原点建立平面直角坐标系，$T$的坐标变为$\left(3, -\dfrac{3}{2}\right)$，椭圆方程变为$\dfrac{(x+1)^2}{4}+\dfrac{\left(y+\dfrac{3}{2}\right)^2}{3} = 1$，即$3x^2+6x+4y^2+12y=0$

                设$PQ$方程为$mx+ny=1$，根据$PQ$过$T$可知满足$3m-\dfrac{3}{2}n = 1$

                利用$PQ$方程将椭圆方程齐次化有$3x^2+6x(mx+ny)+4y^2+12y(mx+ny)=0$，即$(3+6m)x^2+(4+12n)y^2+(6n+12m)xy = 0$

                同除得到$(3+6m)+(4+12n)k^2+(6n+12m)k = 0$

                根据韦达定理可得$k_1+k_2 = \dfrac{-(6n+12m)}{4+12n}$，代入$3m-\dfrac{3}{2}n = 1$可得$k_1+k_2 = -1$
            \end{dingautolist}
        \end{enumerate}
    \end{jieda}
    \item 已知椭圆 $\Gamma  : \dfrac{{x}^{2}}{{a}^{2}} + \dfrac{{y}^{2}}{5} = 1\left( {a > \sqrt{5}}\right) ,M\left( {0,m}\right) \left( {m > 0}\right) ,A$ 是 $\Gamma$ 的右顶点.
    \begin{enumerate}
        \item 若 $\Gamma$ 的焦点$(2,0)$,求离心率 $e$
        \item 若 $a = 4$ ,且 $\Gamma$ 上存在一点 $P$ ,满足 $\vv{PA} = 2\vv{MP}$ ,求 $m$
        \item 已知 ${AM}$ 的中垂线 $l$ 的斜率为 $2,l$ 与 $\Gamma$ 交于 $C\text{、}D$ 两点, $\angle {CMD}$ 为钝角,求 $a$ 的取值范围.
    \end{enumerate}
    \begin{jieda}
        \begin{enumerate}
            \item $c = 2$，$b^2 = 4$，故$a = 3$，所以$e = \dfrac{2}{3}$
            \item $a = 4$则椭圆的方程为$\dfrac{x^2}{16} + \dfrac{y^2}{5} = 1$，$A(4, 0)$，设$P(x_0, y_0)$，则$\left\{\begin{array}{l}
                 \vv{PA} = (4-x_0, -y_0)  \\
                 \vv{MP} = (x_0, y_0-m)
            \end{array}\right.$，因为$\vv{PA} = 2\vv{MP}$，所以$\left\{\begin{array}{l}
                 4-x_0 = 2x_0  \\
                 -y_0 = 2(y_0-m)
            \end{array}\right.$，即$\left\{\begin{array}{l}
                 x_0 = \dfrac{4}{3}  \\
                 m = \dfrac{3y_0}{2}
            \end{array}\right.$，将$x_0 = \dfrac{4}{3}$代入椭圆方程，可得$y_0 = \pm \dfrac{2}{3}\sqrt{10}$，故$m = \sqrt{10}$（负舍）
            \item 因为${AM}$ 的中垂线 $l$ 的斜率为 $2$，所以$k_{AM} = -\dfrac{1}{2}$，因此$m = \dfrac{a}{2}$，故中垂线方程为$y = {2x} - \dfrac{3}{4}a$ ,显然直线 $l$ 过椭圆内点 $\left( {\dfrac{3}{8}a,0}\right)$ ,
                故直线与椭圆恒有两不同交点,
                设 $C\left( {{x}_{1},{y}_{1}}\right) ,D\left( {{x}_{2},{y}_{2}}\right)$ ,则$\left\{\begin{array}{l}
                     \vv{MC} = \left(x_1, y_1 - \dfrac{a}{2}\right)  \\
                     \vv{MD} = \left(x_2, y_2 - \dfrac{a}{2}\right)
                \end{array}\right.$
                因为 $\angle {CMD}$ 为钝角,则 $\vv{MC} \cdot  \vv{MD} < 0$ ,且 $M\left( {0,\dfrac{a}{2}}\right)$ ,
                则有 ${x}_{1}{x}_{2} + \left( {{y}_{1} - \dfrac{a}{2}}\right) \left( {{y}_{2} - \dfrac{a}{2}}\right)  < 0$ ,
                所以 ${x}_{1}{x}_{2} + \left( {2{x}_{1} - \dfrac{5a}{4}}\right) \left( {2{x}_{2} - \dfrac{5a}{4}}\right)  = 5{x}_{1}{x}_{2} - \dfrac{5}{2}a\left( {{x}_{1} + {x}_{2}}\right)  + \dfrac{25}{16}{a}^{2} < 0$ 
                
                联立$l$和椭圆方程有$\left\{  \begin{array}{l} y = {2x} - \dfrac{3}{4}a \\  5{x}^{2} + {a}^{2}{y}^{2} = 5{a}^{2} \end{array}\right.$ 消 $y$ 得 $\left( {4{a}^{2} + 5}\right) {x}^{2} - 3{a}^{3}x + \dfrac{9}{16}{a}^{4} - 5{a}^{2} = 0$ ,
                由韦达定理得 $\left \{\begin{array}{l}
                     {x}_{1} + {x}_{2} = \dfrac{3{a}^{3}}{4{a}^{2} + 5} \\
                     {x}_{1}{x}_{2} = \dfrac{\dfrac{9}{16}{a}^{4} - 5{a}^{2}}{4{a}^{2} + 5}
                \end{array}\right .$，代入可得$5\left( {\dfrac{9}{16}{a}^{4} - 5{a}^{2}}\right)  - \dfrac{5a}{2} \times  3{a}^{3} + \dfrac{25}{16}{a}^{2}\left( {4{a}^{2} + 5}\right)  < 0$ ,解得 ${a}^{2} < {11}$ ,
                又 $a > \sqrt{5}$ ,
                故 $\sqrt{5} < a < \sqrt{11}$ ,即 $a$ 的取值范围是 $\left( {\sqrt{5},\sqrt{11}}\right)$ .
        \end{enumerate}
    \end{jieda}

        \item 已知双曲线 $\Gamma  : {x}^{2} - \dfrac{{y}^{2}}{{b}^{2}} = 1\left( {b > 0}\right)$ ,左、右顶点分别为 ${A}_{1},{A}_{2}$ ,过点 $M\left( {-2,0}\right)$ 的直线交双曲线 $\Gamma$ 于 $P,Q$ 两点.
    \begin{enumerate}
        \item 若 $\Gamma$ 的离心率为 2 ,求 $b$
        \item 若 $b = \dfrac{2\sqrt{6}}{3}$ , ${\Delta M}{A}_{2}P$ 为等腰三角形,且点 $P$ 在第一象限,求点 $P$ 的坐标
        \item 连接 ${QO}$ ( $O$ 为坐标原点)并延长交 $\Gamma$ 于点 $R$ ,若 $\vv{{A}_{1}R} \cdot  \vv{{A}_{2}P} = 1$ ,求 $b$ 的最大值
    \end{enumerate}
    \begin{jieda}
        \begin{enumerate}
            \item 由双曲线的方程知 $a = 1,c = \sqrt{1 + {b}^{2}}$，因为离心率为2,所以 $\dfrac{c}{a} = \dfrac{\sqrt{1 + {b}^{2}}}{1} = 2$ ,得 $b = \sqrt{3}$ .
            \item 当 $b = \dfrac{2\sqrt{6}}{3}$ 时,双曲线 $\Gamma  : {x}^{2} - \dfrac{3{y}^{2}}{8} = 1$ ,且 ${A}_{2}\left( {1,0}\right)$ .
            因为点 $P$ 在第一象限,所以 $\angle P{A}_{2}M$ 为钝角.
            又 ${\triangle M}{A}_{2}P$ 为等腰三角形,所以 $\left| {{A}_{2}P}\right|  = \left| {{A}_{2}M}\right|  = 3$ .
            设点 $P\left( {{x}_{0},{y}_{0}}\right)$ ,且 ${x}_{0} > 0,{y}_{0} > 0$ ,则 $\left\{  \begin{matrix} \sqrt{{\left( {x}_{0} - 1\right) }^{2} + {y}_{0}^{2}} = 3, \\  {x}_{0}^{2} - \dfrac{3{y}_{0}^{2}}{8} = 1, \end{matrix}\right.$
            得 $\left\{  \begin{matrix} {x}_{0} = 2 \\  {y}_{0} = 2\sqrt{2} \end{matrix}\right.$ ,所以 $P\left( {2,2\sqrt{2}}\right)$ .
            
            \item 由双曲线的方程知 ${A}_{1}\left( {-1,0}\right) ,{A}_{2}\left( {1,0}\right)$ ,且由题意知 $Q,R$ 关于原点对称.
            
            设 $P\left( {{x}_{1},{y}_{1}}\right) ,Q\left( {{x}_{2},{y}_{2}}\right)$ ,则 $R\left( {-{x}_{2}, - {y}_{2}}\right)$. 易知直线 ${PQ}$ 与 $y$ 轴垂直时不符合题意，故可设直线 ${PQ}$ 的方程为 $x = {my} - 2$ .

            $\vv{{A}_{1}R} = \left( {-{x}_{2} + 1, - {y}_{2}}\right) ,\vv{{A}_{2}P} = \left( {{x}_{1} - 1,{y}_{1}}\right) .$
            由 $\vv{{A}_{1}R} \cdot  \vv{{A}_{2}P} = 1$ ,得 $\left( {-{x}_{2} + 1}\right) \left( {{x}_{1} - 1}\right)  - {y}_{1}{y}_{2} = 1$ ,
            所以 $\left( {{x}_{2} - 1}\right) \left( {{x}_{1} - 1}\right)  + {y}_{1}{y}_{2} =  - 1$ ,即 $\left( {m{y}_{2} - 3}\right) \left( {m{y}_{1} - 3}\right)  + {y}_{1}{y}_{2} =  - 1$，整理得 $\left(m^2 + 1\right) {y}_{1}{y}_{2} - {3m}\left( {{y}_{1} + {y}_{2}}\right)  + {10} = 0$

            联立直线与双曲线的方程得 $\left\{  \begin{array}{l} x = {my} - 2, \\  {x}^{2} - \dfrac{{y}^{2}}{{b}^{2}} = 1, \end{array}\right.$
            消去 $x$ ,得 $\left( {{b}^{2}{m}^{2} - 1}\right) {y}^{2} - 4{b}^{2}{my} + 3{b}^{2} = 0$ ,
            且 ${b}^{2}{m}^{2} - 1 \neq  0$ ,即 ${m}^{2} \neq  \dfrac{1}{{b}^{2}}$ ,根据韦达定理有$\left\{\begin{array}{l}
                 {y}_{1} + {y}_{2} = \dfrac{4{b}^{2}m}{{b}^{2}{m}^{2} - 1}  \\
                 {y}_{1}{y}_{2} = \dfrac{3{b}^{2}}{{b}^{2}{m}^{2} - 1} 
            \end{array}\right.$ .
            
            代入方程$\left(m^2 + 1\right) {y}_{1}{y}_{2} - {3m}\left( {{y}_{1} + {y}_{2}}\right)  + {10} = 0$，得$\left( {{m}^{2} + 1}\right)  \cdot  \dfrac{3{b}^{2}}{{b}^{2}{m}^{2} - 1} - {3m} \cdot  \dfrac{4{b}^{2}m}{{b}^{2}{m}^{2} - 1} + {10} = 0$ ,
            整理得 ${b}^{2}m^2 + 3{b}^{2} - {10} = 0$ ,所以 ${b}^{2} = \dfrac{10}{{m}^{2} + 3} \leq \dfrac{10}{3}$，取最大值时$b = \dfrac{\sqrt{30}}{3}, m = 0$，满足要求，故 $b$ 的最大值为 $\dfrac{\sqrt{30}}{3}$ .
        \end{enumerate}
    \end{jieda}
    
\end{enumerate}